% Options for packages loaded elsewhere
\PassOptionsToPackage{unicode}{hyperref}
\PassOptionsToPackage{hyphens}{url}
%
\documentclass[
]{article}
\usepackage{amsmath,amssymb}
\usepackage{iftex}
\ifPDFTeX
  \usepackage[T1]{fontenc}
  \usepackage[utf8]{inputenc}
  \usepackage{textcomp} % provide euro and other symbols
\else % if luatex or xetex
  \usepackage{unicode-math} % this also loads fontspec
  \defaultfontfeatures{Scale=MatchLowercase}
  \defaultfontfeatures[\rmfamily]{Ligatures=TeX,Scale=1}
\fi
\usepackage{lmodern}
\ifPDFTeX\else
  % xetex/luatex font selection
\fi
% Use upquote if available, for straight quotes in verbatim environments
\IfFileExists{upquote.sty}{\usepackage{upquote}}{}
\IfFileExists{microtype.sty}{% use microtype if available
  \usepackage[]{microtype}
  \UseMicrotypeSet[protrusion]{basicmath} % disable protrusion for tt fonts
}{}
\makeatletter
\@ifundefined{KOMAClassName}{% if non-KOMA class
  \IfFileExists{parskip.sty}{%
    \usepackage{parskip}
  }{% else
    \setlength{\parindent}{0pt}
    \setlength{\parskip}{6pt plus 2pt minus 1pt}}
}{% if KOMA class
  \KOMAoptions{parskip=half}}
\makeatother
\usepackage{xcolor}
\setlength{\emergencystretch}{3em} % prevent overfull lines
\providecommand{\tightlist}{%
  \setlength{\itemsep}{0pt}\setlength{\parskip}{0pt}}
\setcounter{secnumdepth}{-\maxdimen} % remove section numbering
\usepackage{amsmath,amssymb}
\usepackage{microtype}
\usepackage{geometry}
\geometry{margin=1in}
\setcounter{secnumdepth}{0}
\setlength{\emergencystretch}{3em}
\ifLuaTeX
  \usepackage{selnolig}  % disable illegal ligatures
\fi
\usepackage{bookmark}
\IfFileExists{xurl.sty}{\usepackage{xurl}}{} % add URL line breaks if available
\urlstyle{same}
\hypersetup{
  pdftitle={Path-independent unnormalized energy at zero volume},
  hidelinks,
  pdfcreator={LaTeX via pandoc}}

\title{Path-independent unnormalized energy at zero volume}
\author{}
\date{4 October 2026}

\begin{document}
\maketitle

\begin{abstract}
For every pseudoeffective zero-volume class on a smooth compact connected
K\"ahler manifold, we construct a path-independent unnormalized energy
after choosing an actual reference singularity type and its energy
origin. The actual reference current determines a degree \(d<n\).
A common relative full-mass domain and a genuine finite-energy domain
are independent of the auxiliary K\"ahler form. The energy extends from
bounded comparison type to this whole domain by finite mixed integrals
and relative truncations. We prove its cocycle, monotonicity, concavity,
decreasing continuity, gauge invariance and exact anchor transformation.
A finite-dimensional quotient of complementary cohomology powers is
the necessary and sufficient minimum auxiliary-reference datum for the
entire energy family. Constants retain positive raw mass. The bounded
perturbation limit and the stated iterated finite-energy limit recover
this construction; a concentration example rules out an unrestricted
joint moving-form limit.
\end{abstract}

\textbf{The question and the construction.} Original P15.1 asks for a
path-independent unnormalized energy with minimum necessary reference
scale and independence of its stated perturbation choices. Here its
carrier is \(X\), of complex dimension \(n\). Given
\(P=\theta+dd^c\psi\), put \(m=n-d\). The direct construction is
\[
 E_{\psi,q(\lambda^m)}(u)=\frac1{d+1}\sum_{a=0}^d
 \int_X(u-\psi)\langle T_u^aP^{d-a}\rangle\omega^m,
 \qquad \lambda=[\omega],\quad u\in\mathcal F^1_\psi.
\]
The common domain, finite integrals and minimum quotient are constructed
below. The anchor chooses the energy origin; changing it uses the proved
cocycle. The denominator in the perturbation limit is
\(\binom nd\varepsilon^m\). The leading actual reference mass term is
\(\binom nd M_d\varepsilon^m\), so the construction retains \(M_d\)
without probability normalization. Genuinely unbounded inputs are
possible, as a later example proves; they need not occur for every
reference type.

\section{n1. Foundations and reconstruction of the bounded
prescribed-type
primitive}\label{n1.-foundations-and-reconstruction-of-the-bounded-prescribed-type-primitive}

This paper addresses P15.1 with precisely retained actual singularity
data and approximation hypotheses. The bounded construction is
reconstructed because it is indispensable to the common-domain and
reference results in n2 and n3. The proof uses the explicitly stated
published foundations; no earlier draft or review verdict is a
mathematical hypothesis.

Let X be a connected smooth compact Kähler manifold of complex dimension
n\textgreater=1. In formulas we use \[
d^c=\frac{i}{4\pi}(\bar\partial-\partial),\qquad dd^c=\frac{i}{2\pi}\partial\bar\partial.
\] Fix a pseudoeffective class alpha with volume zero, a smooth closed
real representative theta, and an actual theta-psh psi not identically
minus infinity. Put P=T\_psi and T\_u=theta+dd\^{}c u. The retained type
is the two-sided, extended-valued comparison \[
\mathcal C_\psi=\{u\in\operatorname{PSH}(X,\theta):\psi-C\leq u\leq\psi+C
\text{ for some }C<\infty\}.
\] The polar loci of these potentials agree. Differences on polar sets
can be set arbitrarily when integrated against nonpluripolar measures.
The empty current product is 1.

\subsection{Exact external foundations and
applicability}\label{exact-external-foundations-and-applicability}

\textbf{B (nonpluripolar products).}
Boucksom--Eyssidieux--Guedj--Zeriahi, \emph{Monge--Ampère equations in
big cohomology classes}, Acta Math. 205 (2010), 199--262,
arXiv:0812.3674v2, Definition 1.1, Propositions 1.4 and 1.6, and Theorem
1.8: on compact Kähler X the nonpluripolar product of k\textless=n
closed positive (1,1) currents is finite, positive and closed,
symmetric, multilinear under sums of positive currents, and plurifinely
local. Its coefficient measures, equivalently its trace measures, do not
charge pluripolar sets. Smooth positive factors are ordinary wedge
factors. These statements concern actual currents. All factors used
below are closed positive actual currents on compact Kähler X. No
multilinearity of the positive product of classes is inferred.

\textbf{D (scalar mixed mass theorems).} Darvas--Di Nezza--Lu,
\emph{Monotonicity of nonpluripolar products and complex Monge--Ampère
equations with prescribed singularity}, Anal. PDE 11 (2018), 2049--2087,
Proposition 2.1 and Theorem 2.4: if theta\_j are n smooth closed real
(1,1) forms with pseudoeffective classes and u\_j,v\_j are theta\_j-psh,
equal bounded singularity types in every slot imply equal integrals of
the top mixed products. If u\_j is less singular than v\_j in every
slot, the u-product integral is at least the v-product integral. These
statements explicitly include pseudoeffective classes and impose neither
bigness nor a small-unbounded-locus condition. In each application to k
current slots the other n-k slots are the indicated smooth Kähler forms
with zero potentials; those classes are Kähler and hence
pseudoeffective. DDNL's operator is i partial bar-partial; potentials
u/(2 pi) represent the same currents in that convention. Positive
rescaling preserves all singularity comparisons.

\textbf{V (boundary volume).} BEGZ Sections 1.3--1.4 and Definition
1.17: a big class admits minimal-singularity potentials, less singular
than any potential in that class; its volume is their top nonpluripolar
mass. For pseudoeffective alpha the positive product and volume are the
decreasing limits along addition of any Kähler class. Thus
vol(alpha+epsilon lambda) tends to vol(alpha) as epsilon decreases to
zero. For nef classes the positive product is the ordinary cup product.
D bounds the mass of any actual positive current in a big class by the
minimal-singularity mass. We use only this bound and the boundary limit,
never equality of arbitrary actual mass with class volume.

\textbf{I (bounded admissible integration by parts).} D.-V. Vu,
\emph{Convexity of the class of currents with finite relative energy},
arXiv:2005.13241v3 (25 March 2022), Section 2, Lemmas 2.5--2.6, Theorem
2.7, with Proposition 2.1(iii): if R is closed positive of bidegree
(n-1,n-1), and a,b are bounded R-admissible dsh functions, then the
finite intrinsic Dirichlet products are defined, the diagonal one is
positive, and \[
\int b\langle dd^ca\mathbin{\dot\wedge}R\rangle
=\int a\langle dd^cb\mathbin{\dot\wedge}R\rangle
=-\int\langle db\wedge d^ca\mathbin{\dot\wedge}R\rangle.\tag{I}
\] Here a dsh function is a difference of quasi-psh functions, and
R-admissibility means that its two representation poles have zero R
trace mass. For R formed from ordinary nonpluripolar factors times
smooth positive factors, iterated relative products give exactly the
signed difference of the two ordinary products with the varying slot
inserted. A smooth R also has no polar mass. Accordingly all bounded
differences and smooth test functions in our uses are admissible. Smooth
functions are dsh by adding a sufficiently large common Kähler
background to their Hessians. Finiteness of signed products follows from
B. No unbounded IBP is used.

Vu uses dd\^{}c\_V=i partial bar-partial/pi, twice our operator, as
directly checked on the supplied page-4 image of arXiv:2004.11111v1.
Source potentials u/2 represent our currents. Source gradient d(a/2)
wedge d\^{}c\_V(a/2) is one half of our gradient; multiply source IBP by
2 to obtain (I). The intrinsic square identity is \[
da\wedge d^ca\wedge R=\tfrac12dd^c(a^2)\wedge R-a\,dd^ca\wedge R.\tag{G}
\] The source v3 Equation (2.3) has this factor, whereas its repeated
page-6 and page-11 formulas omit it. Those defects remain literal in the
annex. We use (G). For local nonnegative bounded psh A,B representing
a=A-B the correct square coefficients are 1,1,-1/2: \[
dd^c A^2R+dd^c B^2R-\tfrac12dd^c(A+B)^2R-a(dd^cA-dd^cB)R.
\] For completeness the zero-cutoff application in source IBP still
follows with these coefficients. Take negative representations
phi\_1,\ldots,phi\_4 of a,b, dd\^{}c phi\_j\textgreater=-Omega, and \[
\zeta_k=1+k^{-1}\max(\textstyle\sum_j\phi_j,-k),\quad
a_k=\max(\phi_1,-k)-\max(\phi_2,-k),
\] and similarly b\_k. The nonzero locus of zeta\_k, not necessarily its
closed support, lies in all four strict finite truncation loci. The
differences a\_k,b\_k are uniformly bounded. Bounded Bedford--Taylor
calculus with the intrinsic half-factor and distributional Stokes gives
\[
0\leq\int d\zeta_k\wedge d^c\zeta_k R
=-\int\zeta_k dd^c\zeta_k R
\leq \frac4k\int\Omega R,
\qquad
\int da_k\wedge d^ca_k R\leq2\|a_k\|_\infty\int\Omega R.
\] The first inequality uses dd\^{}c zeta\_k+4 Omega/k\textgreater=0 and
0\textless=zeta\_k\textless=1. The second uses dd\^{}c
max(phi\_j,-k)+Omega\textgreater=0 and the equal cohomological masses of
those two bounded products. Cauchy--Schwarz kills the IBP terms
containing d zeta\_k. On its nonzero locus locality identifies the
intrinsic gradients and signed products with the truncations. The
admissible gradient limits are finite by the cited lemmas; their poles,
and the poles of all positive relative insertions, are uncharged.
Dominated convergence and polarization give (I). This explains the
precise repair needed in our applications; it does not silently alter
the printed sources. The v1 Proposition 4.2 index and half-factor
defects, and the printed Cauchy--Schwarz defect in v1 Theorem 3.7, are
retained but those proofs are not used here.

\textbf{H (smooth dd\^{}c lemma).} Deligne--Griffiths--Morgan--Sullivan,
\emph{Real homotopy theory of Kähler manifolds}, Invent. Math. 29
(1975), 245--274, Lemma 5.11 and Remark 5.14: a smooth form closed under
d and d\^{}c and d-exact is dd\^{}c-exact on a compact Kähler manifold,
also over the reals. Applied to a smooth real closed cohomologically
zero (1,1) form, its primitive has degree zero; the (1,1) type and real
part give a smooth real function, with the operator scalar adjusted.
These are precisely our hypotheses when comparing representatives.
Complete selected proof and preceding Kähler identities from the
supplied Stony Brook author-hosted copy are in the annex; no
source-status verdict is assumed.

\textbf{F (relative full mass and weighted convexity).} Vu v3 Section 3
defines E\_P(R) for tuples (S\_1,\ldots,S\_k) in the respective
reference classes, more singular than (P\_1,\ldots,P\_k), by equality of
cohomology classes of their R-relative nonpluripolar products. For chi
convex increasing on R with chi(-infinity)=-infinity, the joint
chi-energy is the sum over all subsets J of \[
\int -\chi\!\left(\sum_{j=1}^k(u_j-\varphi_j)\right)
\left\langle\prod_{j\in J}S_j\prod_{j\notin J}P_j\mathbin{\dot\wedge}R\right\rangle.
\] Theorem 3.4 says that the diagonal full-mass and diagonal weighted
classes are convex. Corollary 3.2 says the weighted class is unchanged
under a cohomologous bounded-equivalent change of the common reference
P. The printed proofs refer to prior arguments; we use these as
precisely scoped literature theorems, not as newly reconstructed
results. Our applications have k=d\textgreater=1, all P\_j=P, smooth
closed positive R=omega\^{}(n-d), and chi(t)=t. The product is
top-degree. On connected compact X its cohomology equality is exactly
scalar mass equality because integration identifies H\^{}(2n)(X,R) with
R. Source differences are halved, so source joint linear energy is \[
\frac d2\sum_{a=0}^d\binom da
\int(\psi-u)\langle T_u^aP^{d-a}\rangle\omega^{n-d}
\] after normalizing u\textless=psi and source potentials negative. The
positive factor does not affect finiteness. We never apply this
k\textgreater=1 source statement with k=0.

\textbf{C (capacity and bounded weighted convergence).} Guedj--Zeriahi,
\emph{Intrinsic capacities on compact Kähler manifolds}, J. Geom. Anal.
15 (2005), 607--639: for a Kähler background Omega, Proposition 3.6
bounds the capacity of a negative Omega-psh sublevel \{f\textless-A\} by
a constant times (\textbar\textbar f\textbar\textbar\_(L1)+n int
Omega\^{}n)/A. Proposition 3.7 gives capacity convergence for decreasing
\textbf{bounded} Omega-psh potentials with bounded limit. Corollary 3.8
gives quasicontinuity; Proposition 3.10 compares local Bedford--Taylor
and global Kähler capacities on a finite relatively compact chart cover.
DDNL Proposition 2.2 states that uniformly bounded local psh factors
converging in capacity in all n slots, with uniformly bounded
quasicontinuous weights converging in capacity, have weakly convergent
weighted Bedford--Taylor top products. All local applications below
first truncate every product factor and insert the fixed smooth Kähler
factors to reach n slots.

We spell out the general decreasing-capacity deduction needed later. If
f\_j decreases to a nontrivial quasi-psh f in a common smooth
lower-Hessian background, choose a dominating Kähler Omega and subtract
a common upper-bound constant. The bounded truncations max(f\_j,-A)
decrease to max(f,-A), so Proposition 3.7 applies. Outside
\{f\textless=-A\}, truncations and originals agree. Proposition 3.6,
with for example \{f\textless-A+1\}, bounds the discarded set by O(1/A).
First send j to infinity and then A to infinity. This proves general
decreasing capacity convergence. It is a deduction using bounded
Proposition 3.7 and the sublevel estimate, not the false attribution of
unbounded hypotheses to that proposition. Proposition 3.10 and Corollary
3.8 justify the local capacity/quasicontinuity uses.

\textbf{L (comparison-set locality with correct scope).} Vu v3 Theorem
3.3: if product factors are bounded local psh functions, equal on
W=intersection \{v\_j\textgreater v'\_j\}, where the comparison
potentials v\_j,v'\_j may be unbounded psh, their products relative to a
fixed positive current agree on W. To use this with unbounded
nonpluripolar factors below, first truncate every local product
potential at -k. On W intersected with all original and replacement
product potentials \textgreater-k, the bounded factors agree. Theorem
3.3 applies, adding these strict comparisons to W. B's defining
finite-locus restrictions identify the bounded products with the
nonpluripolar ones there. Exhaust the simultaneous finite loci; the
remaining poles have zero mass. This proves exactly the comparison-set
locality we use, without granting an unbounded-factor version of Theorem
3.3. For d=0 locality is simply equality of the fixed smooth measure.

\subsection{Degree, scalar masses and mixed
vanishing}\label{degree-scalar-masses-and-mixed-vanishing}

For a Kähler form omega define \[
M_k(\psi,[\omega])=\int\langle P^k\rangle\omega^{n-k},\quad
d=\max\{k:M_k>0\},\quad m=n-d.
\] M\_0=int omega\^{}n\textgreater0. If omega' is Kähler, c
omega\textless=omega'\textless=C omega for positive constants.
Telescoping the differences of the positive powers gives \[
c^{n-k}M_k(\psi,[\omega])\leq M_k(\psi,[\omega'])
\leq C^{n-k}M_k(\psi,[\omega]).
\] This applies to every positive product current; hence d is
independent of the form. Since P+epsilon omega\textgreater=epsilon omega
its class is big, and B,D,V give \[
0\leq\sum_{k=0}^n\binom nk\varepsilon^{n-k}M_k
=\int\langle(P+\varepsilon\omega)^n\rangle
\leq\operatorname{vol}(\alpha+\varepsilon[\omega])\longrightarrow0.
\] Thus M\_n=0 and d\textless n.~D, filled with n-k smooth slots, says
that every k-fold mixed product from C\_psi has trace mass M\_k. If the
factors are merely more singular than psi, its mass is at most M\_k. For
k\textgreater d this proves vanishing as a current: a positive current
with zero strictly positive trace mass has all coefficient measures
zero, since their absolute values are locally bounded by a fixed
multiple of its trace. For bounded-equivalent anchors all M\_k and d
agree by D. Closedness from B and distributional Stokes imply that M\_k
depends only on the auxiliary class and is a homogeneous polynomial of
degree n-k. None of these degrees is asserted to be class numerical
dimension.

\subsection{Bounded energy, cocycle and representative independence:
full
reconstruction}\label{bounded-energy-cocycle-and-representative-independence-full-reconstruction}

Fix a smooth complementary product W=omega\_1\ldots omega\_(n-q) of
Kähler forms and 0\textless=q\textless=n.~On the convex hull of finitely
many bounded-type potentials u\_0,\ldots,u\_N, put h\_i=u\_i-u\_0 and
u\_t=u\_0+sum t\_i h\_i. The coefficient integrals in \[
A_i(t)=\int h_i\langle T_{u_t}^q\rangle W
\] are finite because h\_i is bounded and B gives finite uncharged
measures. Multilinearity makes A\_i a polynomial on the simplex. For
q\textgreater=1, signed differentiation in a current slot and I, with
R=\textless T\_(u\_t)\^{}(q-1)\textgreater W of bidegree (n-1,n-1), give
\[
\partial_j A_i=q\int h_i\langle dd^ch_j\mathbin{\dot\wedge}R\rangle
=q\int h_j\langle dd^ch_i\mathbin{\dot\wedge}R\rangle
=\partial_i A_j.
\] R is smooth if q=1 and otherwise nonpluripolar, so the bounded
differences are admissible. For q=0 all A\_i are constant. Ordinary
Stokes for this closed polynomial one-form on each finite triangle gives
the cocycle for its segment integrals. Polynomial extension and
continuity include simplex boundaries. The binomial expansion and beta
integral give the segment value \[
E^q_{v,W}(u)=\int_0^1\int(u-v)
\langle((1-t)T_v+tT_u)^q\rangle W\,dt
=\frac1{q+1}\sum_{a=0}^q\int(u-v)\langle T_u^aT_v^{q-a}\rangle W.\tag{1}
\] Indeed \(\binom qa\int_0^1 t^a(1-t)^{q-a}dt=1/(q+1)\). At q=d and
W=omega\^{}m write E\_(psi,lambda)(u), lambda={[}omega{]}. The triangle
identity proves \[
E_{\psi,\lambda}(u)-E_{\psi,\lambda}(v)=E_{v,\lambda}(u).\tag{2}
\] Along a segment \(u_t=u+t h\) the derivative is
\(\int h\langle T_{u_t}^d\rangle\omega^m\). Its second derivative for
d\textgreater=1 is \[
-d\int\langle dh\wedge d^ch\mathbin{\dot\wedge}
(\langle T_{u_t}^{d-1}\rangle\omega^m)\rangle\leq0.
\] This uses the positive intrinsic diagonal in I. For d=0 energy is
affine. Positivity of the first derivative for nonnegative directions
proves monotonicity. D supplies every same-type mixed degree-d mass
M\_d, so constants give c M\_d and pointwise sandwiches give the M\_d
Lipschitz estimate in bounded differences. These derivative claims are
only for segments and finite-dimensional hulls.

If one Kähler factor in W=omega\_1\ldots omega\_m changes by dd\^{}c g
within its class, H supplies smooth real g. For each
Q=\textless T\_u\^{}a P\^{}(d-a)\textgreater{} set R=Q times the other
m-1 factors. Its bidegree is (n-1,n-1), it is closed positive and has no
polar mass. The bounded h=u-psi and g are admissible. I gives \[
\int h\,dd^cg\,R
=\int g\langle dd^ch\mathbin{\dot\wedge}R\rangle=0.\tag{3}
\] The signed insertion is
\((\langle T_u^{a+1}P^{d-a}\rangle-\langle T_u^a P^{d-a+1}\rangle)\)
times the other factors. Both (d+1)-fold products vanish. This also
works at d=0, where R is smooth and both one-fold nonpluripolar products
vanish. Thus the bounded energy integral with m mixed Kähler factors is
representative independent in every slot. It is symmetric multilinear in
their classes. Decompose an arbitrary smooth closed real (1,1) form eta
as (eta+A Omega)-A Omega, both Kähler for large A. This extends the
integral multilinearly to arbitrary factors. Changing a representative
inserts dd\^{}c g; expanding the remaining differences of Kähler factors
reduces to (3). Hence the extension is well-defined, independent of
temporary decompositions, representatives and bases. Its diagonal is a
homogeneous degree-m polynomial in V=H\^{}(1,1)(X,R).

Simultaneously replacing theta by theta+dd\^{}c g and (u,psi) by
(u-g,psi-g) preserves all currents and differences. This proves gauge
invariance; a bounded-equivalent anchor change has (2), not absolute
anchor independence.

\subsection{Checked bounded perturbation
prescription}\label{checked-bounded-perturbation-prescription}

For u in C\_psi and epsilon\textgreater0, the relative unnormalized top
energy in theta+epsilon omega has the exact expansion \[
E^{n,\varepsilon,\omega}_\psi(u)
=\sum_{k=0}^d\binom nk\varepsilon^{n-k}E^k_{\psi,\omega^{n-k}}(u).\tag{4}
\] One can expand (1) at degree n along the positive segment; the
coefficient of its k-fold current mixed term is binom(n,k)/(k+1).
Equivalently \[
\sum_{j=0}^n\binom ja\binom{n-j}{k-a}=\binom{n+1}{k+1},
\] by counting a (k+1)-subset of \{0,\ldots,n\} by the location j of its
(a+1)-st element. Invalid binomial indices mean zero. High-degree terms
vanish by the preceding proof. Each lower coefficient is bounded by
\(\|u-\psi\|_\infty M_k\). The actual perturbed anchor mass is
sum\_(k\textless=d) binom(n,k)epsilon\^{}(n-k)M\_k, without equality to
class volume being assumed.

If {[}omega\_epsilon{]} tends to a Kähler lambda, the cohomological M\_k
polynomials are bounded. The degree-d energy polynomial is continuous in
class. Formula (4) therefore proves the leading energy limit after
dividing by binom(n,d)epsilon\^{}m, for this fixed bounded-type u.
Moving perturbed psh potentials of the same type, uniformly within
delta\_epsilon of u with delta\_epsilon-\textgreater0, have energy error
at most delta\_epsilon times the actual anchor mass: apply the
positive-background cocycle and monotonicity already proved at arbitrary
q=n.~This gives the same limit even without representative bounds. No
division by vol(alpha)=0 occurs, and the lower coefficients themselves
need not be representative independent. This bounded statement will not
be extended to uncontrolled moving forms for unbounded potentials.

Support: none. Full source bodies, hashes, inputs and rerun evidence are
supplied in the assembled P15.1-argument.md and local audit artifacts.
This note's analytic claims use the exactly stated external theorems,
not their prior review status.

\section{n2. One full-mass domain and one genuine finite-energy domain
for all auxiliary
forms}\label{n2.-one-full-mass-domain-and-one-genuine-finite-energy-domain-for-all-auxiliary-forms}

This note uses n1, including its exactly scoped literature foundations
B,D,F,C,L,I,H and reconstructed bounded energy. It proves the proposed
common-domain statements and reconstructs the finite-energy extension
required for n3. The actual anchor type remains fixed, with the anchor
transformation subsequently proved by the cocycle.

\subsection{The common full-mass domain: scalar proof in every mixed
slot}\label{the-common-full-mass-domain-scalar-proof-in-every-mixed-slot}

Set m=n-d\textgreater=1 and \[
\mathcal E_\psi(\omega)=\left\{u\in\operatorname{PSH}(X,\theta):u\leq\psi+C,
\ \int\langle T_u^d\rangle\omega^m=\int\langle P^d\rangle\omega^m\right\}.
\] Assume first d\textgreater=1 and that the equality holds for
omega\_0. Given any Kähler omega\_1, compactness allows a constant
A\textgreater0 for which D\_0=A omega\_0-omega\_1 is itself a smooth
Kähler form: choose A above the maximum eigenvalue of omega\_1 relative
to omega\_0. For 0\textless=j\textless=m put \[
\Delta_j=\int\big(\langle P^d\rangle-\langle T_u^d\rangle\big)
\omega_1^j D_0^{m-j}.
\] These are scalar quantities, not claimed positive currents. Apply
DDNL Theorem 2.4 with n slots: d copies of background theta and
potentials psi versus u; j copies of background omega\_1, with potential
zero on both sides; and m-j copies of background D\_0, also with
potential zero. Every background class is pseudoeffective, each
potential is psh in that background, and u\textless=psi+C means psi is
less singular in the varying slots. The smooth slots have equal types.
Consequently every Delta\_j\textgreater=0. Ordinary smooth expansion
gives \[
0=A^m\int(\langle P^d\rangle-\langle T_u^d\rangle)\omega_0^m
=\sum_{j=0}^m\binom mj\Delta_j.
\] All coefficients are strictly positive; hence every Delta\_j=0, in
particular Delta\_m=0. Reversing the roles gives equality of the
domains. If d=0 both products are 1, each Delta\_j is identically zero,
and full mass is automatic. Thus a single domain E\_psi works for every
smooth Kähler form and every auxiliary Kähler class. This proves the
unverified proposal without presuming positivity of a current
difference.

This also gives mixed complementary equalities. If
omega\_1,\ldots,omega\_m are arbitrary Kähler forms, apply the just
proved equality to omega(t)=sum t\_i omega\_i for all t\_i\textgreater0.
Both sides are polynomials in t. Equality on the nonempty open orthant
implies equality of every coefficient, including the coefficient of
t\_1\ldots t\_m. Therefore \[
\int\langle T_u^d\rangle\omega_1\cdots\omega_m
=\int\langle P^d\rangle\omega_1\cdots\omega_m.\tag{5}
\] Only ordinary smooth-factor polynomial algebra is used here.

\subsection{Common finite-energy domain and exact agreement with joint
linear
energy}\label{common-finite-energy-domain-and-exact-agreement-with-joint-linear-energy}

Define F\^{}1\_psi by u in E\_psi and, for one Kähler omega, \[
\sum_{a=0}^d\binom da\int|u-\psi|
\langle T_u^aP^{d-a}\rangle\omega^m<\infty.\tag{6}
\] For d=0 (6) means int \textbar u-psi\textbar{}
omega\^{}n\textless infinity. Quasi-psh functions are locally integrable
against smooth volume; a finite chart cover gives this integral for
every such u,psi. Thus in this case F\^{}1\_psi=E\_psi, with no
application of Vu's k\textgreater=1 weighted theorem.

For any fixed omega\_0,omega\_1, choose c,C\textgreater0 with c
omega\_0\textless=omega\_1\textless=C omega\_0. Telescoping positive
powers gives inequalities of complementary positive forms, and wedging
each positive mixed d-current yields inequalities of measures \[
c^m\langle T_u^aP^{d-a}\rangle\omega_0^m
\leq\langle T_u^aP^{d-a}\rangle\omega_1^m
\leq C^m\langle T_u^aP^{d-a}\rangle\omega_0^m.
\] Integrate the nonnegative weight \textbar u-psi\textbar. Finiteness
of each term is therefore equivalent for the two forms; common full mass
was proved above. Hence F\^{}1\_psi is genuinely independent of the
auxiliary class and representative, not merely independent within one
class. A product of any m fixed Kähler forms is bounded above and below
by positive multiples of omega\_0\^{}m by the same factorwise
comparison, so all the corresponding mixed complementary weighted
integrals are finite as well.

Constants cause no problem: they do not change currents, and every mixed
mass is finite, bounded by the reference mass by D. Thus (6) is
invariant under shifting u or psi by finite constants. Normalize
u\textless=psi and both source potentials negative. For d\textgreater=1,
specialization of Vu's Definition (3.1) with k=d, diagonal S=T\_u,
common P and linear chi(t)=t gives exactly \[
\frac d2\sum_{a=0}^d\binom da\int(\psi-u)
\langle T_u^aP^{d-a}\rangle\omega^m.
\] The factor 1/2 is the checked source-potential convention. All
summands are nonnegative and their finitely many coefficients are
positive, so finiteness is exactly (6), also exactly the unweighted
sum's finiteness. Source full mass is top-cohomology equality and hence
exactly the scalar equality in E\_psi. The source's diagonal joint
linear class and our F\^{}1\_psi therefore agree precisely after
normalization. Vu Theorem 3.4 applies with smooth R=omega\^{}m and
chi(t)=t; it gives convexity of both common domains. For d=0 convexity
follows directly from the more-singular comparison and smooth-volume
integrability.

\subsection{Mixed full mass and mixed weighted integrability on the
whole
domain}\label{mixed-full-mass-and-mixed-weighted-integrability-on-the-whole-domain}

For a finite family u\_1,\ldots,u\_N in E\_psi, subtract constants to
make each \textless=psi and put w=(u\_1+\ldots+u\_N)/N. At
d\textgreater=1 source full-mass convexity makes w full mass. By B, the
mass of \textless T\_w\^{}d\textgreater{} omega\^{}m is the average of
the N\^{}d ordered mixed masses. Each is at most M\_d by DDNL Theorem
2.4, using d theta slots (psi versus the chosen u\_i in each) and m
identical smooth omega slots. They are nonnegative, and their average
equals M\_d.~Each must equal M\_d.~This proves every d-fold mixed mass
from the family, with repetitions, for every omega. Polynomial
extraction as in (5) proves the same equality with any m mixed
complementary Kähler factors. d=0 has the fixed empty current product
and is immediate. Adding back constants does not change these masses. In
particular all mixed reference/current masses used in the following
arguments have exactly the specified raw reference mass.

If the family lies in F\^{}1\_psi, weighted convexity also puts w in
F\^{}1\_psi. The term a=d of its energy says int
(psi-w)\textless T\_w\^{}d\textgreater{} omega\^{}m\textless infinity.
Its positive expansion has finitely many terms, so psi-w is integrable
against every ordered product from that expansion. Since \[
0\leq\psi-u_i\leq N(\psi-w),
\] psi-u\_i is integrable against every d-fold mixed product formed from
the family. In particular \textbar u\_i-u\_j\textbar{} is integrable
against all of them. Mixed smooth complementary factors are controlled
by the preceding pointwise comparison. For d=0 local L1 integrability
against smooth volume proves everything directly. Adding constants
preserves finiteness by the mixed mass equalities. These arguments
include psi and all bounded-type anchors as members of the finite
family. Higher-degree mixed products of more-singular factors vanish by
n1's checked scalar DDNL application; no lower-degree weighted
integrability is concluded.

\subsection{Exact truncation identity and the finite-energy
meaning}\label{exact-truncation-identity-and-the-finite-energy-meaning}

For normalized u\textless=psi in E\_psi, set z=psi-u\textgreater=0 off
the poles, u\_L=max(u,psi-L), L\textgreater0, and \[
\mu_a=\langle T_u^aP^{d-a}\rangle\omega^m,\qquad
\mu_{a,L}=\langle T_{u_L}^aP^{d-a}\rangle\omega^m.
\] The maximum is theta-psh by the local maximum property;
psi-L\textless=u\_L\textless=psi, so u\_L is in C\_psi. Both measures
have mass M\_d by the preceding mixed full-mass proof (or by D's
sandwich between u and psi). On the strict comparison set
G\_L=\{u\textgreater psi-L\}, they agree by n1's fully checked
finite-locus use of L and B. Explicitly local potentials rho+u and
rho+psi-L are psh comparison potentials; product factors are first
replaced by bounded local truncations, the simultaneous finite loci are
exhausted, and only then are the nonpluripolar restrictions identified.
Poles have zero mass, including when d=0. No unbounded product factors
are put into the printed bounded-factor theorem.

On B\_L=X minus G\_L, outside the uncharged poles, u\_L-psi=-L. Equality
of total masses and equality on G\_L imply mu\_(a,L)(B\_L)=mu\_a(B\_L).
Thus even without finite energy, \[
\int(\psi-u_L)d\mu_{a,L}
=\int_{G_L}z\,d\mu_a+L\mu_a(B_L)
=\int\min(z,L)d\mu_a.\tag{7}
\] The boundary \{u=psi-L\} is included in B\_L and its weight is
exactly L; locality is used only on the strict set. The argument works
verbatim with a fixed mixed complementary product W and its common mass
from (5).

For u in F\^{}1\_psi define the finite real energy \[
E_{\psi,\lambda}(u)=\frac1{d+1}\sum_{a=0}^d
\int(u-\psi)\langle T_u^aP^{d-a}\rangle\omega^m,\quad\omega\in\lambda.\tag{8}
\] Monotone convergence in (7) proves \[
E_{\psi,\lambda}(u_L)\downarrow E_{\psi,\lambda}(u).\tag{9}
\] It also proves the precise converse: on the full-mass domain,
finiteness of the decreasing bounded truncation-energy limit is
equivalent to (6). Indeed the negative of that limit is the sum of the
nonnegative integrals int z dmu\_a divided by d+1; a finite sum is
finite exactly when every term is finite, equivalently when its positive
binomial-weighted sum is finite. For d=0 this is the ordinary smooth L1
formula. Thus this is a genuine F\^{}1 domain defined equally by a
finite energy value, not an arbitrary renaming of a weighted class.

Monotonicity follows by ordering the truncations and using the bounded
result. Every mixed mass in (8) is M\_d, so \[
E_{\psi,\lambda}(u+c)=E_{\psi,\lambda}(u)+cM_d.\tag{10}
\] For a general u\textless=psi+C use v=u-C: u\_L=v\_(L+C)+C for large
L. Equation (10) and normalized (9) give (9) for general u. The same
argument applies to mixed complementary W. No choice of real truncation
levels tending to infinity changes the limit.

\subsection{Bounded decreasing continuity with an unbounded reference:
complete tail
proof}\label{bounded-decreasing-continuity-with-an-unbounded-reference-complete-tail-proof}

Suppose v\_j decreases to v and all are uniformly bounded relative to
psi: \textbar v\_j-psi\textbar,\textbar v-psi\textbar\textless=B. For a
fixed degree-q product formed from v\_j and psi and fixed n-q smooth
Kähler factors, denote the positive measures by nu\_j,nu. D gives
exactly the same total mass b for all of them. Shift psi to be negative
and put \[
\eta_A=\min(1,\max(\psi+A,0)).
\] The nonzero locus \{psi\textgreater-A\} contains all places where
this cutoff is used. Its closed support may meet the boundary; we do not
claim otherwise. On a relatively compact chart with theta=dd\^{}c rho,
every local psh rho+v\_j,rho+v,rho+psi on this nonzero locus has a
uniform lower bound depending on A,B,rho, and a uniform upper bound
because v\_j\textless=v\_1. Truncate all these local psh factors below a
still lower constant. The resulting factors are uniformly bounded psh,
converge in capacity by n1's C deduction, and give the same products
wherever eta\_A is nonzero. The fixed smooth factors have bounded smooth
local psh potentials as well.

The cutoff can be expressed using the bounded truncation of psi and
smooth rho. Similarly eta\_A(v\_j-psi) equals eta\_A times the
difference of the bounded local psh truncations just chosen. Such
weights are uniformly bounded quasicontinuous functions converging in
capacity: max, min and products on bounded ranges preserve capacity
convergence by their elementary Lipschitz estimates and a union bound on
exceptional sets. A smooth partition of unity is an additional fixed
bounded continuous weight. DDNL Proposition 2.2 therefore proves local,
and then global, convergence of \[
\int\eta_A d\nu_j\ \text{ and }\ \int\eta_A(v_j-\psi)d\nu_j
\] to their limiting counterparts for every fixed A. Bounded-factor
hypotheses hold in every product slot. This is not a direct application
of a big-class convergence theorem in the pseudoeffective background.

The tail weight is bounded by B. Hence its absolute integral is at most
\[
B\left(b-\int\eta_A d\nu_j\right).
\] First let j tend to infinity at fixed A; then let A tend to infinity.
The limiting tail mass tends to zero because nu gives no mass to psi's
polar set. This proves convergence of every full weighted integral and
hence bounded decreasing energy continuity. For q=0 the same conclusion
follows by dominated convergence against smooth volume; the cutoff
argument also applies with no varying product factors.

\subsection{Decreasing continuity, concavity and the unbounded
cocycle}\label{decreasing-continuity-concavity-and-the-unbounded-cocycle}

Let v\_j decrease to u with v\_j,u in F\^{}1\_psi. They have one upper
bound psi+C because v\_1 does. Monotonicity gives \[
E(u)\leq E(v_j)\leq E(\max(v_j,\psi-L)).
\] For fixed L the right-hand sequence decreases to u\_L within a
uniform bounded type, so the lemma proves limsup
E(v\_j)\textless=E(u\_L). Sending L to infinity using (9) proves
E(v\_j)-\textgreater E(u). This includes every decreasing bounded-type
approximation with limit in F\^{}1\_psi. If w=(1-t)u+t v, source
weighted convexity puts w in the domain. The bounded-type sequence
(1-t)u\_L+t v\_L decreases to w; bounded concavity, (9) and decreasing
continuity prove concavity on the entire domain. At d=0 it is directly
affine.

For normalized u,v\textless=psi take simultaneous truncations,
nu\_a=\textless T\_u\^{}a T\_v\^{}(d-a)\textgreater omega\^{}m and
nu\_(a,L)=\textless T\_(u\_L)\^{}a
T\_(v\_L)\^{}(d-a)\textgreater omega\^{}m. They have mass M\_d and agree
on G\_L=\{u\textgreater psi-L\} intersection \{v\textgreater psi-L\} by
the finite-locus comparison proof. Their complementary masses therefore
agree. On B\_L their truncated weight
\textbar u\_L-v\_L\textbar\textless=L, and with z\_u=psi-u,z\_v=psi-v,
\[
L\nu_{a,L}(B_L)=L\nu_a(B_L)
\leq\int_{B_L}(z_u+z_v)d\nu_a\longrightarrow0.
\] The last limit is dominated convergence using the proved mixed
weighted integrability; B\_L decreases to uncharged polar sets. The
original \textbar u-v\textbar{} tail also tends to zero. On G\_L
products and weights agree. Thus \[
\frac1{d+1}\sum_a\int(u_L-v_L)d\nu_{a,L}
\longrightarrow\frac1{d+1}\sum_a\int(u-v)d\nu_a.
\] Taking limits in the reconstructed bounded cocycle yields the finite
unbounded formula \[
E_{\psi,\lambda}(u)-E_{\psi,\lambda}(v)
=E_{v,\lambda}(u)
:=\frac1{d+1}\sum_{a=0}^d\int(u-v)\langle T_u^aT_v^{d-a}\rangle\omega^m.\tag{11}
\] Constants reduce the general case to this normalization: products do
not change and every mixed mass is M\_d.~For d=0 this is the linear
smooth-volume identity.

\subsection{Anchors, representatives and theta
gauge}\label{anchors-representatives-and-theta-gauge}

Let chi in C\_psi. The more-singular relations to psi and chi are
identical up to constants, and bounded mixed mass equality gives their
identical reference mass and degree. Thus their full-mass domains agree
for all forms. Every chi belongs to F\^{}1\_psi because its differences
are bounded. For u in F\^{}1\_psi, the mixed-family integrability proof
applied to (u,chi,psi) shows that
\textbar u-chi\textbar\textless=\textbar u-psi\textbar+\textbar\textbar psi-chi\textbar\textbar\_infinity
is integrable against every product \textless T\_u\^{}a
T\_chi\^{}(d-a)\textgreater omega\^{}m. Hence u is in F\^{}1\_chi.
Applying the same argument with the roles reversed proves equality of
finite-energy domains. This is a direct anchor-domain proof, consistent
with Vu Corollary 3.2, and does not depend on an unbounded IBP
assertion.

Taking v=chi in (11) gives exactly \[
E_{\chi,\lambda}(u)=E_{\psi,\lambda}(u)-E_{\psi,\lambda}(\chi).\tag{12}
\] The mixed-integral definition with anchor chi is the displayed E\_v;
all terms have just been proved finite. At d=0 smooth L1 integrability
proves the same facts without source weighted theorems.

For cohomologous Kähler representatives omega,omega', the domain is
already common. Each u\_L bounded energy is representative independent
by n1. Equation (9), valid for both forms, proves representative
independence of the whole finite-energy value. Simultaneous smooth theta
gauge changes leave actual currents, relative differences, comparisons
and truncations invariant, so all constructions are unchanged.

\subsection{A genuinely unbounded
member}\label{a-genuinely-unbounded-member}

On \(X=\mathbb P^1\times\mathbb P^1\) take normalized Fubini--Study
\(\kappa\) with \(\int\kappa=1\), \(\theta=p_1^*\kappa\), \(\psi=0\) and
\(\omega=p_1^*\kappa+p_2^*\kappa\). The smooth semipositive theta makes
its class nef; its square is zero, so B,V give volume zero. The actual
degree is d=1 and M\_1=1. Choose 0\textless a\textless1/2 and in a small
disk at zero set \[
v(z)=-b(-\log|z|)^a
\] near zero. Continue by a smooth cutoff on a fixed annulus. On the
punctured inner disk the radial second derivative with respect to
t=log\textbar z\textbar{} is b a(1-a)(-t)\^{}(a-2)\textgreater0. The
annular negative Hessian is bounded by a constant times b, so small
b\textgreater0 ensures S=kappa+dd\^{}c v\textgreater=0 globally. Define
v(0)=-infinity and subtract a constant to make v\textless=0. The locally
integrable subharmonic extension has no atom at zero: its outward radial
flux is proportional to b a(-log r)\^{}(a-1), tending to zero.
Distributional Stokes gives S(X)=1, and no mass is lost in its
nonpluripolar part. It is smooth off zero.

Smooth-volume integrability follows from exponential area decay in
s=-log r. Its weighted singular-current integral is a positive constant
times \[
\int^\infty s^a s^{a-2}ds=\int^\infty s^{2a-2}ds<\infty.
\] Thus \(u=p_1^*v\) has full degree-one mass and both finite weighted
terms in (6), but u-psi is unbounded. The construction is an example
demonstrating genuine extension; no fibration hypothesis appears in the
general theorem. This analytic proof, not a numerical test, establishes
the example.

Support: n1. These results remove the auxiliary dependence of the
full-mass and finite-energy domains. They supply no unrestricted
unbounded moving-auxiliary perturbation limit.

\section{n3. Coefficientwise extension and the minimum reference for the
entire energy
family}\label{n3.-coefficientwise-extension-and-the-minimum-reference-for-the-entire-energy-family}

This note uses n1 and n2. It proves the coefficientwise unbounded
extension and the exact necessary/sufficient finite-dimensional
reference criterion for the preceding construction. The minimum is for
this entire anchored prescribed-type energy family. It is not an
independent geometric classification.

\subsection{Intrinsic bounded functionals and an unisolvent Kähler
set}\label{intrinsic-bounded-functionals-and-an-unisolvent-kuxe4hler-set}

Let V=H\^{}(1,1)(X,R), with its real Bott--Chern/de Rham interpretation
on compact Kähler X, and B=Sym\^{}m V, m=n-d\textgreater=1. V is
finite-dimensional as a subspace of finite-dimensional de Rham
cohomology; H identifies cohomologically equal smooth (1,1)
representatives by a smooth dd\^{}c gauge. For u in C\_psi, define the
symmetric functional ell\_u\^{}psi in B\^{}* by \[
\ell_u^\psi(\lambda_1\odot\cdots\odot\lambda_m)
=\frac1{d+1}\sum_{a=0}^d\int(u-\psi)
\langle T_u^aP^{d-a}\rangle\eta_1\cdots\eta_m,\qquad[\eta_i]=\lambda_i.\tag{13}
\] Here eta\_i are arbitrary smooth closed real representatives. The
extension by differences of Kähler forms and the representative-change
proof in n1 prove well-definedness in all slots; no temporary basis or
chosen decomposition is data. It is a symmetric m-linear form and hence
defines exactly a linear functional on Sym\^{}m V. Its diagonal
polynomial p\_u(lambda)=ell\_u\^{}psi(lambda\^{}odot m) equals
E\_(psi,lambda)(u) on the Kähler cone.

Let h=dim V and N=dim B=binom(h+m-1,m). The tensors lambda\^{}odot m
with lambda Kähler span B. To prove this, if their span were proper,
finite-dimensional duality would give a nonzero f in B\^{}* annihilating
them all. Its associated homogeneous polynomial f(lambda\^{}odot m)
would vanish on the nonempty open Kähler cone. A real polynomial
vanishing on an open set is zero: in coordinates choose a small open
rectangular box in that set, regard it as a polynomial in the first
coordinate, and use the one-variable statement that a nonzero polynomial
cannot have infinitely many roots; its coefficients vanish on the
remaining box. Induct on the number of coordinates. Thus all polynomial
coefficients vanish, and polarization forces f=0, a contradiction. This
last implication is also explicit in the monomial basis: coefficient of
x\^{}beta is (m!/beta!)f(e\_1\^{}odot beta\_1 \ldots{} e\_h\^{}odot
beta\_h).

Consequently there exist N Kähler classes lambda\_1,\ldots,lambda\_N
whose powers form a basis of B: iteratively choose a power outside the
span of those already chosen. The evaluation map \[
\mathcal A:B^*\longrightarrow\mathbb R^N,\quad
f\longmapsto(f(\lambda_i^{\odot m}))_{i=1}^N\tag{14}
\] is a linear isomorphism. This is the promised unisolvent set. Its
existence uses only openness and finite-dimensional polynomial algebra.
Choosing a different such set changes coordinates, not the intrinsic
functional.

\subsection{Coefficientwise convergence on the common unbounded
domain}\label{coefficientwise-convergence-on-the-common-unbounded-domain}

Fix u in F\^{}1\_psi and u\_L=max(u,psi-L). n2 proves pointwise in every
Kähler class \[
\ell_{u_L}^\psi(\lambda^{\odot m})
=E_{\psi,\lambda}(u_L)\longrightarrow E_{\psi,\lambda}(u)\in\mathbb R.
\] In particular the N evaluations (14) converge. Apply the fixed
inverse A\^{}(-1). It follows that ell\_(u\_L)\^{}psi converges in
B\^{}*, equivalently coefficientwise in any finite basis, to one
ell\_u\^{}psi. Evaluating the limit at any Kähler lambda and using the
same n2 limit proves \[
E_{\psi,\lambda}(u)=\ell_u^\psi(\lambda^{\odot m}).\tag{15}
\] Thus the entire unbounded finite-energy value is a homogeneous
degree-m class polynomial. No nonuniform lower-epsilon coefficients
enter this argument. For any other unisolvent set or basis, the
resulting functional has the same diagonal values on the Kähler cone;
the open-set polynomial proof gives uniqueness. n2 also gives the same
pointwise limit for every decreasing bounded-type approximation with
limit u, so all such approximations give exactly this coefficient limit.
Arbitrary real truncation levels tending to infinity are covered by (9).

The functionals have the explicit integral formula (13) on the unbounded
domain as well. Here is a proof that also checks convergence of the
mixed integrals. For a mixed product W of m Kähler representatives, n2
gives the full mass of every product \textless T\_u\^{}a
P\^{}(d-a)\textgreater W and its truncated counterpart, and their
locality on G\_L. Thus (7) holds with omega\^{}m replaced by W. Each
weight is integrable by n2, so each weighted integral for u\_L converges
to the integral for u. For arbitrary smooth real representatives eta\_i,
choose A\_i large enough that eta\_i+A\_i Omega and A\_i Omega are
Kähler. Expand their product into its 2\^{}m signed positive products W.
Every weight integral is finite because W is bounded by a constant times
Omega\^{}m, and all 2\^{}m truncation limits just proved exist. The
bounded integral functionals converge to ell\_u\^{}psi coefficientwise.
Taking those integral limits gives exactly (13). This also proves
representative and decomposition independence for these explicit
unbounded mixed integrals without using unbounded IBP. Equivalently they
are obtained from (15) by the intrinsic polarization formula \[
\ell_u^\psi(\lambda_1\odot\cdots\odot\lambda_m)
=\frac1{m!}\sum_{J\subset\{1,\ldots,m\}}(-1)^{m-|J|}
p_u\!\left(\sum_{j\in J}\lambda_j\right),\tag{16}
\] where p\_u is the unique polynomial on all V. The empty sum is 0.
Formula (16) follows by expanding the polynomial and observing that the
inclusion-exclusion sum kills every term missing a variable, leaving m!
times the fully multilinear term. No meaning as a positive reference is
assigned to a non-Kähler argument of this polynomial extension.

\subsection{No extra reference directions appear at finite
energy}\label{no-extra-reference-directions-appear-at-finite-energy}

Define intrinsically \[
S_\psi=\operatorname{span}\{\ell_u^\psi:u\in\mathcal C_\psi\}\subset B^*,
\quad K_\psi=\operatorname{Ann}(S_\psi)\subset B,
\quad Q_\psi=B/K_\psi,
\] with quotient map q. Every ell\_(u\_L)\^{}psi belongs to S\_psi. This
subspace is finite-dimensional and closed, so its coefficient limit
ell\_u\^{}psi belongs to S\_psi for every u in F\^{}1\_psi. Conversely
C\_psi is contained in F\^{}1\_psi by bounded masses and bounded
weights. Therefore \[
\operatorname{span}_{u\in\mathcal F^1_\psi}\ell_u^\psi=S_\psi,
\qquad
\bigcap_{u\in\mathcal F^1_\psi}\ker\ell_u^\psi=K_\psi.\tag{17}
\] This equality is the substantive extension result, beyond defining a
quotient on bounded potentials.

For two Kähler classes lambda\_0,lambda\_1, equations (15)--(17) give
the exact equivalence \[
\big[E_{\psi,\lambda_0}(u)=E_{\psi,\lambda_1}(u)
\text{ for every }u\in\mathcal F^1_\psi\big]
\quad\Longleftrightarrow\quad
q(\lambda_0^{\odot m})=q(\lambda_1^{\odot m}).\tag{18}
\] For the forward direction restrict to bounded u; their functionals
span S, so the tensor difference lies in K. For the backward direction
every unbounded functional belongs to S and annihilates K. Hence the
quotient criterion is both necessary and sufficient on the entire
genuine domain.

\subsection{Raw mass, constants, anchor change and gauge
change}\label{raw-mass-constants-anchor-change-and-gauge-change}

Define the mass functional in B\^{}* by \[
\mathfrak M_\psi(\lambda_1\odot\cdots\odot\lambda_m)
=\int\langle P^d\rangle\eta_1\cdots\eta_m.
\] Closedness of \textless P\^{}d\textgreater{} and ordinary smooth
Stokes prove representative independence. Its diagonal
M\_d(psi,lambda)\textgreater0 on the Kähler cone, and its mixed Kähler
values are positive because the strictly positive complementary product
dominates a positive multiple of a Kähler trace power of the nonzero
positive current \textless P\^{}d\textgreater. At d=0 the current is 1
and the same statement is the positive smooth volume integral. The mass
functional is retained in S: \[
\ell_{\psi+c}^\psi=c\mathfrak M_\psi,\qquad
\ell_{u+c}^\psi=\ell_u^\psi+c\mathfrak M_\psi.\tag{19}
\] Indeed currents are unchanged by constants, and every degree-d mixed
mass from (13) is exactly the reference mixed mass by n2. This proves
(19) initially on mixed Kähler arguments and then on B by
spanning/polarization. In particular psi+1 belongs to C\_psi and
supplies mathfrak M\_psi. Thus Q retains the positive raw mass and a
scale cannot be eliminated by forgetting constants. K is contained in
ker mathfrak M\_psi and Q has dimension at least one. There is no
normalization of M\_d or of zero volume to a probability.

Let chi in C\_psi. n2 gives the identical domains and the scalar anchor
cocycle for every Kähler lambda and every u in the domain. Unique
polarization gives the functional identity on B \[
\ell_u^\chi=\ell_u^\psi-\ell_\chi^\psi.\tag{20}
\] Every bounded ell\_u\^{}chi lies in S\_psi. Conversely
ell\_psi\^{}chi=-ell\_chi\^{}psi and
ell\_u\^{}psi=ell\_u\^{}chi-ell\_psi\^{}chi, so S\_chi=S\_psi. Hence
K\_chi=K\_psi and the quotient, its admissible image, and the minimum
equivalence are anchor invariant within the bounded type. The origin of
energy still requires an anchor, changing by (20). Mixed bounded mass
equality shows mathfrak M\_chi=mathfrak M\_psi. This proves all anchor
statements including d=0, without asserting that C\_psi in that case
consists only of constant translates.

A simultaneous theta gauge change theta'=theta+dd\^{}c g, u'=u-g,
psi'=psi-g, with smooth real g, preserves every actual current,
difference, truncation and integral in (13). It transports C\_psi and
F\^{}1\_psi bijectively to the new potential domains, preserves all
functionals and identifies the same S,K,Q. Any two smooth real
representatives of alpha differ in this way by H. Auxiliary
representative independence was proved above by bounded reconstruction
and truncation, not by a gauge applied to the auxiliary paths. Adding
constants to g does not change the identification of anchored
differences. Keeping theta as a convenient coordinate representative is
permissible; its smooth gauge is not independently necessary reference
data.

\subsection{Exact universal minimum sense and admissible
topology}\label{exact-universal-minimum-sense-and-admissible-topology}

Each ell\_u descends to a linear functional bar-ell\_u on Q. The map \[
\mathcal A_{\rm all}:B\longrightarrow\mathbb R^{\mathcal F^1_\psi},
\quad t\longmapsto(\ell_u(t))_u
\] has kernel K, by (17). It induces an injective linear map from Q onto
this family's finite-dimensional output space. If r:B-\textgreater R\_0
is any linear reference descriptor and for every u there is a linear
f\_u on r(B) with ell\_u=f\_u composed with r, then ker r is contained
in K. Thus there is one well-defined surjective linear map \[
F:r(B)\longrightarrow Q,\qquad F(r(t))=q(t).
\] Well-definedness is exactly ker r subset K; surjectivity follows from
surjectivity of q, and uniqueness follows because r(B) is spanned by
those images. This proves the coarsest linear descriptor property. It
does not assert a geometric classification of K.

For actual auxiliary choices, define the admissible image \[
\mathcal I_\psi=\{q(\lambda^{\odot m}):\lambda\text{ Kähler}\}\subset Q.
\] Every energy is determined by this descriptor. Among arbitrary
set-valued descriptors r\_0 on the Kähler cone, if equality of
r\_0(lambda)'s implies equality of all energy values, (18) yields a
uniquely defined map from r\_0's image onto I\_psi sending r\_0(lambda)
to q(lambda\^{}odot m). Conversely the latter descriptor suffices. This
is minimum informational data for this family, not a proof that this
image is a manifold, a geometric carrier, or all of Q. The map
lambda-\textgreater q(lambda\^{}odot m) need not be injective, and
I\_psi need not be a vector space.

Give Q its finite-dimensional quotient topology and I\_psi the subspace
topology. For fixed u every bar-ell\_u is continuous; hence
tau\_j-\textgreater tau in I implies convergence of its energy values
even when the representing classes do not converge. Conversely choose
bounded u\_1,\ldots,u\_s whose ell's are a basis of S (s=dim Q). Then \[
Q\longrightarrow\mathbb R^s,\quad
\tau\longmapsto(\bar\ell_{u_i}(\tau))_i
\] is a linear isomorphism: its kernel is zero by definition of K, and
both sides have dimension s. Thus convergence of these finitely many
bounded energy coordinates is equivalent to quotient convergence. This
provides the weakest finite-dimensional topology making the whole energy
family continuous.

If u\_j decreases to u with all in F\^{}1\_psi, n2 gives convergence at
every unisolvent lambda\_i; (14) then proves
ell\_(u\_j)-\textgreater ell\_u in B\^{}*. Accordingly their polynomials
converge uniformly on every compact subset of V (bound the evaluation
tensors there), and their descended functionals converge uniformly on
bounded subsets of Q. If simultaneously tau\_j-\textgreater tau in I,
the values bar-ell\_(u\_j)(tau\_j) tend to bar-ell\_u(tau), since
convergent sequences of functionals and points are bounded and the
evaluation pairing is continuous. These are admissible
energy-polynomial/quotient statements. They make no assertion about a
rescaled top energy with epsilon-dependent rough representatives, nor
about variable unbounded integration weights outside the proved
decreasing theory.

\subsection{Scope and a nontrivial quotient
computation}\label{scope-and-a-nontrivial-quotient-computation}

The remaining reference data for this energy family are X, its complex
structure and class alpha, the actual bounded comparison type (with its
degree d), an anchor specifying the additive origin, and tau in I\_psi.
A smooth theta is only a gauge representative. The current integrals are
carried by the actual n-dimensional X with product bidegree (d,d). A
finite-dimensional quotient of reference tensors is not a geometric
d-dimensional carrier. Actual type is not reduced to ideal equivalence
or class numerical dimension.

For clarity the whole auxiliary class need not be minimal. In n2's
\(X=\mathbb P^1_y\times\mathbb P^1_z\) example with
\(\theta=p_1^*\kappa\) and psi=0, V has basis
\(a=[p_1^*\kappa],b=[p_2^*\kappa]\). One direct computation is the
product cell decomposition of the two spheres: there are two real
2-cells, no odd cells, and the two pullback classes evaluate
independently on the two fibers; hence H\^{}2 has precisely these two
generators, both of type (1,1). A real class A a+B b is Kähler exactly
when A,B\textgreater0: restriction to each fiber proves necessity, and
the product form proves sufficiency. Any theta-psh potential restricts
on each z-fiber to a psh function or the identically-minus-infinity
function. A nontrivial psh function on a compact connected curve is
constant by the maximum principle. Thus all fibers are constant and
u(y,z)=v(y), extended-valuedly; a nontrivial restriction to a y-slice is
kappa-psh and identifies v. For bounded u in C\_0 it is a bounded v.
Wedge calculation gives \[
E_{0,Aa+Bb}(p_1^*v)=\frac B2\int_{\mathbb P^1}v\big((\kappa+dd^cv)+\kappa\big).
\] All bounded ell\_u therefore lie in the span of the functional B, and
constants supply that functional itself. Hence S is exactly
one-dimensional, K is the A-axis, Q is R with q(Aa+Bb)=B, and
I=(0,infinity). By (17) this description holds for the entire genuine
F\^{}1 domain, including n2's unbounded example. This is a concrete
geometric computation in one special case, not a universal geometric
kernel formula.

Support: n1, n2. The new conclusion is the common-domain,
coefficientwise and minimum-family-reference theorem with exact scope,
not an original-atom completion.

\section{n4. Analytic limitation, original residuals, and reproducible
execution
evidence}\label{n4.-analytic-limitation-original-residuals-and-reproducible-execution-evidence}

This note uses n1, n2 and n3. It reconstructs the supplied analytic
concentration obstruction rather than treating the prior proof as a
grant, and separates the substantive new theorem from the unchanged
original P15.1. P15.2 and P15.3 are outside this assignment; no minimum
operator-reference construction or result for those atoms is claimed.

\subsection{Same-class concentration still prevents an unrestricted
unbounded perturbation
theorem}\label{same-class-concentration-still-prevents-an-unrestricted-unbounded-perturbation-theorem}

Use the actual unbounded finite-energy potential \(u=p_1^*v\) proved in
n2 on \(X=\mathbb P^1\times\mathbb P^1\), with \(\theta=p_1^*\kappa\),
psi=0, d=1,m=1, int kappa=1 and \(v=-b(-\log|z|)^a-c\) near zero,
0\textless a\textless1/2, c\textgreater=0. Its actual current is
\(p_1^*S\), \(S=\kappa+dd^c v\) of mass one without a pole atom. All
ambient hypotheses, degree-one full mass and finite joint linear energy
were proved there.

For epsilon\textgreater0 sufficiently small put
rho\_epsilon=exp(-epsilon\^{}(-2/a)). Choose a smooth nonnegative area
form beta\_epsilon of integral one supported in
\{\textbar z\textbar\textless rho\_epsilon\}; a disk bump rescaled and
normalized gives it and extension by zero is smooth. Every real 2-form
on this curve is closed and of type (1,1). Set \[
\kappa_\varepsilon=\tfrac12\kappa+\tfrac12\beta_\varepsilon,
\qquad \omega_\varepsilon=p_1^*\kappa_\varepsilon+p_2^*\kappa.
\] These are smooth strictly positive Kähler forms. H\^{}2(P\^{}1,R) is
one-dimensional and integration determines its class; hence
{[}kappa\_epsilon{]}={[}kappa{]}. All auxiliary classes
{[}omega\_epsilon{]} equal the single Kähler
\(\lambda=[p_1^*\kappa+p_2^*\kappa]\). The potential is identical for
every parameter, and T\_u+epsilon omega\_epsilon is positive. At each
fixed parameter v is integrable against kappa\_epsilon because it is
locally L1 and this fixed smooth density is bounded, and against S and
kappa by n2. Thus every weighted mixed top integral in the following
formula is finite, with no uniform bound assumed in epsilon.

Write \(A=p_1^*\kappa\), \(B=p_2^*\kappa\),
\(A_\varepsilon=p_1^*\kappa_\varepsilon\) and \(U=p_1^*S\). Products of
two horizontal (1,1) factors vanish off the pole, hence their
nonpluripolar extensions vanish; mixed products below give no pole mass.
Ordinary smooth multiplication, or the defining nonpluripolar
finite-locus restrictions, therefore gives \[
\begin{aligned}
(T_0+\varepsilon\omega_\varepsilon)^2&=2\varepsilon(A+\varepsilon A_\varepsilon)B,\\
\langle(T_u+\varepsilon\omega_\varepsilon)(T_0+\varepsilon\omega_\varepsilon)\rangle
&=\varepsilon(U+A+2\varepsilon A_\varepsilon)B,\\
\langle(T_u+\varepsilon\omega_\varepsilon)^2\rangle
&=2\varepsilon(U+\varepsilon A_\varepsilon)B.
\end{aligned}
\] The actual top masses of the anchor and u are both 2
epsilon(1+epsilon), so this is a full relative top-mass potential for
every positive perturbation. The joint linear-weight integrals of these
three measures are finite at every fixed parameter by the integrability
already checked. This obstruction uses neither lost perturbed mass nor
infinite energy at a fixed parameter.

Directly summing those three finite weighted terms, divided by 3, gives
\[
E^{2,\varepsilon,\omega_\varepsilon}_0(u)
=2\varepsilon E_{0,\lambda}(u)+2\varepsilon^2\int v\kappa_\varepsilon,
\qquad E_{0,\lambda}(u)=\tfrac12\int v(S+\kappa)\in\mathbb R.
\] Consequently \[
\frac{E^{2,\varepsilon,\omega_\varepsilon}_0(u)}{2\varepsilon}
=E_{0,\lambda}(u)+\varepsilon\int v\kappa_\varepsilon.\tag{21}
\] On the bump's support, -v\textgreater=b epsilon\^{}(-2), and
v\textless=0 globally. Therefore int v
kappa\_epsilon\textless=-(b/2)epsilon\^{}(-2), so the last term tends to
minus infinity. For the fixed smooth product omega instead, the same
expansion has the fixed finite integral int v kappa, and the quotient
tends to E\_(0,lambda)(u). These two paths have exactly the same
auxiliary class, exactly the same quotient descriptor from n3, and
exactly the same potential at every parameter.

This proves an analytic counterexample to extending n1's
moving-representative leading-perturbation assertion to all F\^{}1\_psi
without additional controls. It does not contradict n3's continuity on
the energy quotient: that theorem concerns the degree-d energy
polynomial, whereas (21) has an uncontrolled lower-degree weighted
coefficient. A quotient reference, class convergence, and degree-d
finite energy cannot provide a uniform estimate for that extra term. No
positive probability normalization was used anywhere.

\subsection{Actual degree and leading class-volume coefficients must
remain
distinct}\label{actual-degree-and-leading-class-volume-coefficients-must-remain-distinct}

In the same class choose the actual logarithmic reference on the first
factor \[
g(z)=\log\frac{|z|^2}{1+|z|^2},\qquad\psi=p_1^*g.
\] With kappa=dd\^{}c log(1+\textbar z\textbar\^{}2), int kappa=1, the
distributional identity dd\^{}c log\textbar z\textbar\^{}2=delta\_0
follows from the radial flux of log r\^{}2: its Laplacian flux is 4 pi,
and our dd\^{}c is planar Laplacian times area divided by 4 pi. Thus
kappa+dd\^{}c g=delta\_0, globally, and P is the integration current of
the vertical divisor. Its nonpluripolar degree-one product is zero: it
vanishes off the polar divisor and the defining extension puts no mass
on that divisor. Hence d=0 for this actual reference, whereas the smooth
reference has d=1. The nef class has nonzero first cup product and zero
square, so its nef numerical dimension, defined by its largest nonzero
cup-power trace, is one. Actual-current degree is therefore not class
numerical dimension.

For \(\omega=A p_1^*\kappa+B p_2^*\kappa\), A,B\textgreater0,
actual-current multilinearity gives the perturbed divisor-reference mass
\[
\int\langle(P+\varepsilon\omega)^2\rangle
=2AB\varepsilon^2,
\] while the nef perturbed class has volume by its cup product \[
\operatorname{vol}(\alpha+\varepsilon[\omega])
=2B\varepsilon+2AB\varepsilon^2.
\] Both statements follow from n1's exact volume foundation and the
explicitly computed products. The exponents differ. Constant anchor
shifts have actual perturbed energy c times the first mass, so rescaling
by the leading class-volume exponent would discard their positive actual
degree-zero scale. This explicitly checks why class-volume coefficients
cannot be substituted into the scoped theorem without another proof.

\subsection{Completion with the stated data and path hypotheses}\label{what-is-proved-and-what-remains-for-the-original-atom}

The common-domain and reference statements are proved, with the stipulated actual-type degree and
every fixed smooth Kähler auxiliary choice. The full-mass domain is
independent of all auxiliary classes/forms by scalar mixed deficits.
Joint linear integrability is independent by positive pointwise
comparison. On their common genuine F\^{}1 domain the energy has finite
mixed integrals, exact truncation meaning, cocycle, concavity,
decreasing continuity, gauge invariance and same-type anchor
transformation. Its polynomial coefficients are limits of the bounded
coefficients, and the unbounded energies span exactly the same
finite-dimensional S. The exact reference equivalence on the entire
domain is (18); raw mass survives because constants supply the mass
functional. These domain and coefficient-extension statements are the
substantive progress. Defining an annihilator quotient by itself is not
counted as a new geometric theorem.

The construction retains a prescribed actual comparison type and uses an
actual-current degree and scale. Its minimum quotient is a universal
factorization for this specified family, not a theorem that this is the
minimum independently geometric reference for every possible zero-volume
energy theory. The carrier of its integrals is explicitly X of dimension
n with current bidegree (d,d); no geometric d-dimensional carrier or
universal geometric characterization of the kernel is supplied. The
unchanged original question requires these choices to be explicit and
all claimed paths to be justified; it does not require an intrinsic
\(d\)-dimensional carrier or a minimum across every conceivable theory.

In particular a direct all-F\^{}1, class-convergent
moving-representative perturbation prescription is false by (21).
The finite-energy perturbation prescription used here first takes the
proved bounded leading limit at fixed L and then lets L tend to
infinity. It gives (8), independently of the stated bounded
approximations and reference presentations. No equality with every
joint or reversed limit is asserted. General moving unbounded-potential perturbation limits,
sufficient weighted lower-degree control even for all fixed forms in
arbitrary reference types, and infinite-dimensional path/differential
calculus with unbounded velocities have not been proved. Decreasing
continuity is the precise sequence theorem here; no arbitrary increasing
continuity is asserted. The leading coefficient used in n1 is the
actual-reference polynomial's coefficient, not a general leading
class-volume coefficient. No claim on the optional d=0 rigidity question
is needed or made.

\textbf{Completion.} The construction exists for every pseudoeffective
zero-volume class after its actual reference type is selected. The
minimum quotient determines the whole finite-energy family and retains
raw mass; the direct mixed-integral definition and its cocycle establish
path independence. All claimed perturbation choices satisfy the proved
bounded leading limit or the explicitly chosen iterated truncation
prescription. The common domain, carrier, dimension, energy origin and
reference changes are stated. No uncontrolled simultaneous unbounded
limit is used to define the theory.

For a disconnected compact manifold there are finitely many components.
The same construction applies componentwise, retaining each degree,
scale, domain and origin, and the energy is their sum. Independent
componentwise inputs give the product minimum-reference criterion.

\subsection{Verification and literal source evidence}

The complete proof and literal primary-source annex accompany the
paper. The proof was developed through Xean's native Codex path with
\texttt{gpt-6.1-sol} and \texttt{xhigh} reasoning. A separate source-aware
review reconstructed n1--n4 against the unchanged original P15.1
statement and returned PASS. Its complete report, ten checked
foundations and nineteen source passages are preserved. The source
normalization, gradient half factor, cutoff nonzero locus and bounded
product-factor hypotheses are checked explicitly in n1. The historical
author and earlier scoped reviews remain unchanged.

These computational mathematical audits are not formal proof
certificates or literature-priority claims. The actual TeX compilation
and source hashes establish reproducibility, not mathematical truth.

\end{document}
